\chapter{Introdução}\label{chap:intro}

Conforme os fenômenos da mecânica quântica foram melhor descritos e compreendidos, surgiu um novo paradigma para a solução de problemas: a computação quântica. Essa forma de elaborar algoritmos é equivalente à computação clássica em termos de computabilidade \cite{shor1998quantum}, mas pode apresentar vantagens de complexidade computacional para determinados problemas. Por exemplo, dada uma lista de itens sem organização prévia, a computação clássica verifica a presença de um elemento percorrendo a lista até encontrar o alvo. Essa solução possui complexidade temporal $O(n)$, pois o tempo necessário cresce linearmente em relação ao tamanho da lista. Porém, Grover demonstrou que a mecânica quântica permite resolver esse problema com complexidade temporal $O(\sqrt{n})$ \cite{grover}. Nesse caso, uma lista quatro vezes maior exige apenas o dobro do tempo para localizar o elemento desejado. Portanto, o ganho quadrático resulta da redução da complexidade de $O(n)$ para $O(\sqrt{n})$.

Atualmente, porém, essa área de pesquisa encontra um obstáculo: a dificuldade de executar os algoritmos elaborados, afetando significativamente o desenvolvimento de novos algoritmos, já que testá-los torna-se uma tarefa manual e demorada sem um computador para realizá-la. Essa dificuldade existe porque hardware que possibilite a execução das operações fundamentais da computação quântica é caro e dificilmente acessível. Além disso, mesmo quando acessível, hardware quântico moderno apresenta baixa capacidade, permitindo apenas a execução de algoritmos simples sobre exemplos de casos pequenos \cite{golec2024quantum}.

Dessa forma, para auxiliar no desenvolvimento de algoritmos quânticos, busca-se solução na simulação por meio de computadores clássicos para compensar a falta de hardware quântico. Visto que ambas as formas de definir algoritmos são computacionalmente equivalentes, sabe-se que todo algoritmo quântico pode ser executado em uma máquina clássica e vice-versa. Então, basta definir uma forma de converter algoritmos quânticos em algoritmos clássicos e executar essa nova versão em uma máquina clássica, os quais podem ser executados em máquinas mais acessíveis e baratas. Neste trabalho, essa conversão será feita decompondo o circuito em matrizes, cada porta lógica é representada por uma matriz e a execução do algoritmo se dá pelo produto das matrizes \cite{nielsen2010quantum}.

Por sua vez, a simulação quântica também apresenta desafios. Devido à capacidade dos computadores quânticos de sobrepor um número exponencial de estados possíveis em relação à sua capacidade física, a simulação precisa ser capaz de armazenar uma quantia exponencial de informação, tornando o consumo de memória inviável. De forma a mitigar esse problema, diferentes técnicas de simulação surgem, tentando otimizar o consumo de memória e tempo, como a simulação esparsa \cite{jaques2022leveraging} e o método de tableau \cite{aaronson2004improved}.

Este trabalho consiste na implementação de uma dessas técnicas: o uso de diagramas de decisão para simulação de circuitos quânticos. Nessa técnica, o circuito a ser simulado tem suas portas lógicas convertidas em matrizes e a execução do circuito é representada pela multiplicação sucessiva dessas matrizes. Porém, como mencionado previamente, o tamanho da informação a ser armazenada cresce exponencialmente em relação à largura do circuito e esse crescimento é refletido nas matrizes. Por isso o QMDD (Quantum Multiple-Value Decision Diagram) é útil: essa estrutura de dados representa matrizes de forma concisa, explorando a redundância em sua composição \cite{fujita1997multi}. Essa técnica de compactação apresenta ganho variável, pois matrizes com maior redundância são compactadas de maneira mais eficaz que matrizes com menor redundância. Por exemplo, uma matriz $1000 \times 1000$ preenchida exclusivamente pelo valor $1$ pode ser descrita como uma matriz desse tamanho com todas as posições iguais a $1$, em vez de armazenar individualmente um milhão de números $1$. Na matriz apresentada a seguir, porém, todos os valores são diferentes; desse modo, o QMDD não compacta essa estrutura com a mesma eficácia e consome mais memória.

\begin{equation}
  \label{distinctmatrix}
  \begin{bmatrix}
    1 & 2 & \cdots & 1000 \\
    1001 & 1002 & \cdots & 2000 \\
    \vdots & \vdots & \ddots & \vdots \\
    999001 & 999002 & \cdots & 1000000
  \end{bmatrix}
\end{equation}

No caso das matrizes relevantes para a computação quântica, a repetição de valores é muito comum por conta da natureza das operações que levam às suas construções. Por isso, o QMDD se torna uma estrutura eficiente para realizar a simulação de circuitos quânticos. 

A partir dessas premissas, este trabalho implementa os algoritmos necessários à simulação de circuitos quânticos, incluindo a conversão de cada porta em um QMDD equivalente à matriz da porta, o produto entre dois QMDDs e a medição sobre um QMDD. A implementação é integrada à plataforma Ket \cite{evandro}, uma interface desenvolvida para facilitar a descrição de algoritmos quânticos. Dessa forma, por meio do módulo de simulação por QMDD, o usuário simula algoritmos que não exigem matrizes de baixa redundância de maneira eficiente. Nos casos em que o crescimento exponencial do diagrama é inevitável, como em estados com alto entrelaçamento, a simulação permanece possível, embora seja viável apenas para circuitos de porte reduzido.

Além disso, não é intuitivamente claro quais circuitos levarão a uma matriz com alta ou baixa redundância. Então, este trabalho buscará comparar o desempenho de diferentes circuitos de forma a analisar se existem aspectos em comum entre os circuitos que apresentaram ganho significativo.

\section{Objetivos}

O objetivo geral deste trabalho é implementar um método de simulação adicional para o Ket, que poderá ser escolhido pelo usuário para otimizar a simulação de circuitos cujas representações matriciais possuam grande redundância.

Por sua vez, os objetivos específicos são:
\begin{enumerate}[label=O\arabic*]
  \item Compilar e descrever os algoritmos necessários para o QMDD;

  \item Implementar um simulador utilizando os algoritmos do artigo;

  \item Integrar o simulador produzido ao código fonte do Ket;

  \item Identificar em quais cenários o simulador produzido apresenta vantagem em relação aos métodos de simulação já disponíveis na plataforma.
\end{enumerate}

\section{Estrutura do Documento}

No \cref{chap:compq} são introduzidos os conceitos fundamentais da computação quântica e suas diferentes formas de representação. Ao longo dessa introdução, são destacados os aspectos relevantes para a definição dos requisitos que o QMDD deve atender e os motivos pelos quais a simulação quântica não é trivial. Em seguida, são estabelecidas as funcionalidades necessárias para que uma ferramenta seja adequada à simulação de circuitos quânticos.

Com base nesses requisitos, o \cref{chap:qmdd} apresenta uma introdução ao funcionamento dos grafos, de modo a facilitar o entendimento da estrutura do QMDD, baseada em grafos dirigidos. Em seguida, são descritas as características do QMDD e os algoritmos necessários para implementar as operações exigidas pela simulação de circuitos quânticos.

A seguir, o \cref{chap:development} descreve como os algoritmos introduzidos no capítulo anterior foram implementados. Foram apresentados os detalhes técnicos da implementação, como a linguagem de programação utilizada, os aspectos práticos da implementação, como as estruturas de dados utilizadas para aplicação eficiente dos algoritmos e, por fim, a integração do simulador ao Ket. Além disso, foram descritos os testes realizados para garantir que o simulador implementado produza resultados corretos.

Já o \cref{chap:experiments} apresenta os métodos utilizados para validar o desempenho do simulador implementado, comparando seus resultados com os de outros simuladores disponíveis na plataforma Ket. Nesse capítulo, são apresentados os algoritmos quânticos utilizados nas simulações realizadas para a coleta dos dados.

Por fim, o \cref{chap:conclusion} avalia os resultados finais, pontuando objetivos atingidos e conclusões sobre o método de simulação estudado e implementado ao longo deste trabalho. 

\section{Trabalhos Correlatos}

O uso de Diagramas de Decisão para representação de circuitos quânticos foi inicialmente proposto por \cite{miller2006qmdd}. Neste artigo, é apresentado o uso de uma estrutura baseada em diagramas de decisão para a representação concisa de matrizes que apresentam alta taxa de redundância em seus valores. Além disso, são apresentados os algoritmos necessários para manipular essas matrizes da forma necessária para realização da simulação de circuitos quânticos.

Dada a utilidade do QMDD na execução de funções quânticas, \cite{kydros2025tutorial} apresenta um manual de implementação dessa estrutura. Nesse manual, os algoritmos são definidos formalmente por meio de pseudocódigo e acompanhados das estruturas necessárias para garantir que a implementação seja eficiente. Por exemplo, a definição do uso de tabelas hash para identificação de arestas já existentes. 

\section{Metodologia}

Este trabalho apresenta a construção de um simulador a partir de um método de simulação ainda não utilizado pela plataforma Ket. A construção é realizada por meio da compilação de algoritmos existentes na literatura, de modo a atender aos requisitos do Ket para a simulação de circuitos quânticos.

Os resultados da implementação foram analisados de forma qualitativa e quantitativa. As estatísticas de desempenho do simulador implementado foram comparadas com as dos simuladores já disponíveis no Ket, a fim de verificar se existem cenários nos quais o novo simulador apresenta vantagem. Também foram identificados os cenários com ganho significativo e aqueles com desempenho semelhante ou inferior, permitindo a análise das características dos algoritmos que se beneficiam da representação por QMDD.
